\section{Persona-Based Elicitation Interviews}
\label{sec:pa-elicitation}

Each persona is interviewed by two agents that alternate: the \emph{interviewer} asks one question, the \emph{persona agent} answers it in the role of the persona. The interview of a persona ends when the interviewer decides that the plan is sufficiently covered for this persona, or when the maximum number of turns is reached; then the next persona is interviewed. When all personas are done, the workflow moves to the analyst. Each interview is stored as an \emph{interview transcript} artifact, one per persona, whose items are the messages. Each message has a short identifier, shown to the user, that names the persona and the position of the message in the interview (for example, the fourth message of the first persona), so that later requirements can cite it as evidence.

\subsection{Interviewer}
\label{sec:pa-interviewer}

The interviewer is responsible for covering the plan with each persona. It receives the plan, the persona and the latest turns of the interview, and, in a rework round, the issues that must be clarified. It produces one question per turn, building on the previous answers, and decides when the interview of the persona is complete. Asking a single question per turn is a deliberate design choice: the persona can then answer each question fully, and each answer can be cited on its own as evidence.

\subsection{Persona Agent}
\label{sec:pa-persona-agent}

The persona agent is responsible for answering as the persona. It receives the persona, the question and the interview so far, and produces an answer in the first person that is consistent with the goals, constraints, priorities and concerns of the persona. To counter the agreeableness of aligned models (Section~\ref{sec:bg-pretraining}), it is instructed to disagree where the priorities of the persona conflict with the question, rather than accept whatever the interviewer suggests.

\subsection{Context and Cost}
\label{sec:pa-interview-cost}

The two agents receive only the transcript of the persona being interviewed, not the other interviews. Since the transcript only grows at its end, consecutive calls share their prompt prefix, and the provider serves it from its prompt cache at a lower price. Interviews run one persona after the other rather than in parallel, because parallel calls that start at the same moment cannot reuse each other's cache.
